\documentclass[11pt,a4paper]{article}
\usepackage[margin=25mm]{geometry}
\usepackage{amsmath,amssymb,amsthm,hyperref}
\setlength{\emergencystretch}{2em}
\newtheorem{lemma}{Lemma}
\newtheorem{theorem}[lemma]{Theorem}
\newcommand{\gen}[1]{\langle#1\rangle}
\title{Fracture graphs in the strict high-rank range}
\author{Human Compression project: a component of the L019 recovery}
\date{}
\begin{document}\maketitle
Write $G_i=\gen{\rho_j:j\ne i}$ for a maximal omission of a
string C-marking $\rho_0,\ldots,\rho_{r-1}$. A fracture graph
chooses, for every $i$, one transposition of $\rho_i$ whose endpoints
belong to different $G_i$-orbits. We recover the initial fracture
argument needed for the stability theorem, with the degree-six
action qualification made explicit.

\begin{lemma}\label{wreath}
Let $m\ge5$ be odd. A transitive string C-marking generating
the natural $S_{2m}$ cannot have a maximal omission of rank
$m+1$ with the exceptional $C_2\wr S_m$ marking of the complete
transitive half-degree classification.
\end{lemma}
\begin{proof}
Let the omission be $H=G_i$, preserving its $m$ pairs. The acquired
classification gives its actual base marks
$\lambda, v=\lambda+e_1$, where $\lambda$ flips all pairs.
Its other $m-1$ marks $s_1,\ldots,s_{m-1}$ form a complement
with block images the adjacent transpositions of a path.
The lifts can contain a cocycle; we do not replace the action
by its abstract cocycle normalization. In the actual action,
conjugation on the base still permutes coordinates as the
block image does. The noncommutation diagram of $H$ is the
isolated mark $\lambda$ and the path
\[
                        v,s_1,\ldots,s_{m-1}.                 \tag{1}
\]

The diagram of the parent $S_{2m}$ is connected. Indeed a
disconnection, together with IP, would express this group as a
direct product of two nontrivial commuting normal subgroups.
This is impossible for $S_d$, $d\ge5$: every nontrivial normal
subgroup contains $A_d$, or else centralizes $A_d$ and is trivial;
the natural centralizer of $A_d$ is trivial. Thus every adjacent
pair in the parent string fails to commute. Deleting position
$i$ leaves paths of lengths $i$ and $r-i-1$. To obtain the two
components above, the omitted mark $q=\rho_i$ must lie between
$\lambda$ and an endpoint of (1). In particular $i=1$ or
$r-2$, rather than an endpoint of the whole string.

If the adjacent endpoint is $v$, then $q$ centralizes
$B=\gen{s_1,\ldots,s_{m-1}}$. Apply the acquired actual
pair-centralizer lemma: a group transitive on the pairs, with
primitive block image generated by involutions, has its
centralizer inside the pair stabilizer, except for a regular
prime-dihedral action. Here the block image is $S_m$, so that
exception is impossible. Consequently $q$ preserves the pairs.

If the adjacent endpoint is $s_{m-1}$, then $q$ centralizes
$v,s_1,\ldots,s_{m-2}$. Conjugating $v$ by the latter block
path gives all the actual base vectors
\[
                         \lambda+e_j\quad(1\le j<m).
\]
Their subgroup $N$ has exactly the original pairs as its point
orbits. Every coordinate occurs nontrivially: coordinate $m$
occurs in all these vectors, and each coordinate $j<m$ occurs
in one with a different index. Since $q$ centralizes $N$, it
permutes its orbits and again preserves the pairs. In both
cases $\gen{H,q}$ preserves a nontrivial partition, contradicting
its being the natural $S_{2m}$.
\end{proof}

\begin{lemma}\label{small}
Every connected rank-six C-marking of the natural $S_8$ has
all its maximal omissions intransitive. For every rank-six
C-marking of the natural $A_{11}$, the centralizer in $S_{11}$
of either five-mark endpoint omission is trivial.
\end{lemma}
\begin{proof}
These are finite consequences of two acquired complete
classifications, not assumptions about generic rank bounds.
The accepted degree-eight certificate has exactly two ordered
conjugacy classes, the following marking and its reversal:
\[
 (2\ 3),\ (1\ 2)(3\ 4),\ (4\ 5),\ (5\ 6),\ (6\ 7),\ (7\ 8).
                                                               \tag{2}
\]
Every omission in (2) is intransitive. Omitting its first mark
separates $\{1,2\}$ from the other points; omitting its second
fixes $1$; omitting a later mark separates the corresponding
terminal segment of the displayed path. This proves the first
assertion, including the reverse order.

The complete acquired $A_{11}$ census has four rank-six
ordered representatives. They have common entries
\[
\begin{aligned}
t_0&=(8\ 9)(10\ 11),&t_1&=(7\ 8)(9\ 10),\\
t_3&=(3\ 4)(5\ 6),&t_5&=(1\ 2)(3\ 4),
\end{aligned}
\]
with independently chosen entries
\[
\begin{split}
t_2&=(6\ 7)(10\ 11)\ \text{or}\ (6\ 7)(8\ 9),\\
t_4&=(2\ 3)(4\ 5)\ \text{or}\ (1\ 3)(2\ 5).
\end{split}
                                                               \tag{3}
\]
For all eight endpoint omissions, the literal-permutation
certificate finds exactly one commuting permutation, the
identity. Its exhaustive enumeration assigns the image of an
unassigned point, then propagates that choice under each
specified mark. A conflict or a repeated image rejects the
choice. Every unused image is tried at every new orbit, so
every equivariant bijection is included. A final direct
commutation test checks each completion. We reuse the accepted
point-fixing implementation by adjoining one new fixed point
and fixing it: restriction gives exactly the unrestricted
centralizer on the original eleven points. A separate GAP
centralizer computation confirms all eight orders. The same
replay records the twelve omission orbit partitions for the
two literal $S_8$ representatives. Complete input arrays,
extraction pins and both executions are retained.
\end{proof}

\begin{theorem}\label{main}
Let $G\le S_n$ be transitive, with a string C-marking of rank
$r$, and assume $n\ge7$ and $2r>n+2$. Then $G=S_n$ and every
$G_i$ is intransitive. Every choice of one crossing
transposition for each label is a fracture graph and is a
forest. For $n=5,6$ the same conclusion holds after choosing
the ordinary degree-$n$ action in its marked isomorphism class.
\end{theorem}
\begin{proof}
We acquire the complete transitive half-degree bound (L047)
and the exact alternating maxima (L044). A transitive group
different from $A_n,S_n$ cannot have this rank: its exceptional
wreath rank is $n/2+1$, and the exceptional small actions also
lie below the strict threshold. The alternating maxima are
$4,5,6$ for $n=9,10,11$, no C-marking exists for $n=7,8$,
and the maximum is $\lfloor(n-1)/2\rfloor$ for $n\ge12$.
Thus $G=S_n$.

Suppose $G_i$ is transitive. IP makes it proper in $G$ and
makes its inherited $r-1$ marks an independent C-marking.
Since $2(r-1)>n$, the same two classifications leave only
the following possibilities:
\begin{enumerate}
\item $n=2m$, with $m\ge5$ odd, and $G_i=C_2\wr S_m$
has rank $m+1$;
\item $n=8$, $r=6$, and $G_i$ is the degree-eight imprimitive
exception of rank five;
\item $n=11$, $r=7$, and $G_i=A_{11}$ has rank six.
\end{enumerate}
Lemma~\ref{wreath} excludes the first case. In the second,
the parent diagram is connected, so the complete parent
classification in Lemma~\ref{small} applies and contradicts
transitivity of $G_i$. In the third, the omission diagram is
connected because $A_{11}$ is simple. Since the parent diagram
is a path, $i$ is an endpoint. The omitted involution must
centralize the five marks distant from it, contrary to
Lemma~\ref{small}. All omissions are therefore intransitive.

For each $i$, some transposition of $\rho_i$ joins different
$G_i$-orbits; otherwise every generator would preserve these
orbits and $G$ would be intransitive. Choose one. If the chosen
edges had a cycle, deleting its $i$-edge would leave a path
between that edge's endpoints using only labels different
from $i$. Such a path lies in a single $G_i$-orbit, contradicting
the choice. Thus the selected graph is a forest.

For $n=5,6$, the same complete rank bounds force $G=S_n$.
The independent-family bound $r\le n-1$ forces equality.
The acquired complete small extremal classification gives
only the simplex marking up to marked-group isomorphism.
Choose its ordinary adjacent-transposition action. Its
omission orbits are the two intervals separated by the
omitted transposition, and the path itself is a fracture graph.
\end{proof}

\paragraph{Action qualification and attribution.}
The fixed source is Cameron--Fernandes--Leemans,
\emph{The number of string C-groups of high rank},
arXiv:2212.12723v2, Proposition~4.1. The rank classifications,
the original stability claim and the fracture construction
remain attributed to their authors. The recovered argument
supplies the actual pair-centralizer justification, preserves
cocycle twists in the permutation action, and states the small
action qualification. In degree six the retained census also
contains the outer image of the simplex: both its endpoint
omissions are transitive. Thus a literal assertion about every
degree-six action would be false. This does not change the
source's intended counting by marked-group isomorphism, for
which the ordinary simplex is a valid representative.

The pair-centralizer lemma is the generic first lemma of the
accepted L044 component; its subsequent alternating-group
theorem is not used here. The finite premises are the complete
accepted degree-eight rank-six census, the complete $A_{11}$
census, and the complete degree-five and six extremal census.
Only the stated finite consequences are newly checked.
This component proves neither perfect-split existence nor
the stability bijection or its six constants. It grants no
whole-original or open-conjecture credit and adds no Lean or
compression measurement.
\end{document}
